\section{Discussion}
\label{sec:discussion}
\subsection{Distinguishing supervision, forecast readout, and forecasting fidelity}
The matched two-seed aggregate supports a descriptive hierarchy in macro discrimination: adding future-state supervision improved direct prediction, and adding the complete forecast-attention readout further improved mean AUROC and AP. Patient-bootstrap intervals for the AUROC contrasts against \nojepa{} excluded zero, while the additional \fullmodel{}--\direct{} interval included zero. The complete-versus-control AUROC gain is larger relative to its conditional uncertainty, whereas the incremental forecast-readout gain remains uncertain. The incremental comparison also changes readout width and pooling, so it does not isolate a single causal mechanism.

The archived forecast probes offer counter-evidence to a simple explanation based on more accurate simulated futures. At common epoch 20, auxiliary-supervised direct prediction had better copy-normalized latent forecast skill than the full classifier, while the full classifier had the better holdout pCR ranking. A predicted representation can therefore be useful to classification without outperforming another model under its own latent reconstruction criterion. One hypothesis is that classification gradients shape the forecast state toward endpoint-relevant directions, at the expense of fidelity to some aspects of the EMA target. The nonzero pCR-to-predictor gradient in the full model is consistent with that information pathway. However, a training probe in a separately learned representation space cannot establish this hypothesis, anatomical forecasting quality, or a causal treatment-response mechanism.

These findings suggest concrete controls for future work: classifier-width-matched direct readout, stopped-gradient forecast features, copy-state readout, and independent holdout forecast evaluation against suitable references. The present main experiments predict embeddings rather than generating MRI for the classifier. The internal control label \nojepa{} also retains the shared fusion architecture and transition queries; it removes the future-state auxiliary objective rather than every world-model-related module. Those distinctions should remain explicit when interpreting the ablation.

\subsection{What additional visits change}
The complete model's AUROC increased across T0--T2 on the same patients, with a paired T2--T0 difference of +0.0071. Its sensitivity increased while specificity decreased. Later histories may help identify patients likely to achieve pCR, but they also move the operating point toward more positive predictions. The outcome-conditioned score histograms describe those score changes; they do not show that every patient's prediction becomes more accurate. Later scans may reveal treatment response, but selection of complete-history patients and the retrospective assigned-arm scenario limit how these patterns can be translated into an earlier decision pathway.

Treatment inputs deserve particular care. T0 imaging preceded randomization in the source protocol, whereas the model is conditioned on the eventual assigned arm. This is a retrospective conditional response estimate, not a deployable prediction before treatment allocation. The encoding does not capture actual delivered dose, cycle timing, dose reductions, adherence, or later changes. Changing the arm token would not identify an individual treatment effect without appropriate causal assumptions and treatment-exposure data.

\subsection{Classification and utility cannot be inferred from AUROC alone}
The expanded metrics show why a higher AUROC does not imply a better decision rule. The full model improved macro sensitivity and NPV but reduced specificity and PPV. Direct readout retained slightly higher balanced accuracy and MCC. At the 0.5 operating point, no principal model exceeded the all-negative reference in macro raw accuracy, although all had balanced accuracy above 0.65 and AUROC above 0.71. Raw accuracy would therefore reward ignoring the positive class in this cohort and should not be the principal assessment criterion.

Validation-selected Youden thresholds altered this trade-off without uniformly increasing holdout F1. Their behavior is consistent with the difference between optimizing balanced class separation and optimizing F1. The full model also had worse NLL, Brier score, and ECE, particularly at later landmarks. Hence its ranking benefit should not be presented as calibrated risk estimation. A clinically motivated endpoint would require calibration fitted without holdout leakage and an operating point justified by the consequences of false-positive and false-negative decisions.

Decision curves reinforce that conclusion: auxiliary-supervised direct readout had greater net benefit than the complete model at some hypothetical thresholds, despite its lower AUROC, and the macro net benefit at 0.5 was non-positive for all principal models. These are mathematical trade-off summaries, not evaluations of a defined clinical intervention. They cannot establish a safe de-escalation strategy or support changing care. External calibration and a prespecified action framework are needed before clinical-utility interpretation.

\subsection{Heterogeneity and the limits of subgroup inference}
Biological subgroup results differed in both ranking and threshold behavior. The HR-positive/HER2-negative group showed a positive macro AUROC difference but still low sensitivity at 0.5. Its low pCR prevalence provides context for predictive values, but does not by itself explain the observed sensitivity. The HR-positive/HER2-positive group had a negative full-versus-control AUROC difference. This negative estimate is relevant evidence against a universal benefit and remains in the report. The small HR-negative/HER2-positive group had too little information for precise conclusions. A subgroup whose interval excludes zero and another whose interval includes zero do not, by themselves, demonstrate a significant interaction.

The source-namespace summaries provide a descriptive acquisition-related check within the same public clinical release; they do not create external validation. Crop sensitivity analysis is similarly limited: only one positive holdout outcome occurred among the 11 patients with partial clipping of the largest predicted baseline component, so a reliable AUROC could not be estimated for that group. The appropriate response is to show the denominator and suppress unstable discrimination rather than report a seemingly precise score.

\subsection{Limitations and a confirmatory study}
\label{sec:limitations}
The seeds 21/44 reporting policy was chosen after results were available. Averaging the same pair across models improves consistency of presentation but does not make selection prospective or exhaust training-randomness uncertainty. Patient intervals condition on the fitted models and, for operating-point analyses, fixed validation-selected thresholds. They do not account for retraining, seed or model selection, threshold-selection uncertainty, or the many exploratory subgroup comparisons. The retrospective two-seed study should therefore be followed by a locked analysis on independent patients.

The holdout is an internal downstream split rather than an independent external institution or prospective cohort. Twenty-nine of its patients had belonged to a development split in an earlier lineage. Current downstream supervised models use the new patient-disjoint folds and are not warm-started from those earlier pCR models, but the pretrained codec and localization model have not yet been fully audited for patient overlap with the current holdout. This prevents a claim that every stage of the pipeline has been evaluated on completely unseen patients.

Fixed T0 localization, registration, and crop coverage are further constraints. Geometry and array checks do not establish anatomical registration accuracy. Some baseline predicted components extend beyond the fixed field of view, and the model may omit multifocal disease or changes outside the crop. Holdout eligibility requires complete T0--T2 imaging, while development retains incomplete histories. Missingness, acquisition variability, and subtype prevalence may therefore affect transportability. The training-probe forecast diagnostics do not replace best-checkpoint holdout mechanism evaluation, and their long-horizon denominators are small.

Finally, the empirical comparison remains an internal architectural ablation rather than a full benchmark against clinical-only, conventional imaging, and independent longitudinal methods. Precise pCR pathology criteria, patient-level scan intervals, ethics determination, and upstream provenance still require final documentation. A stronger study would add matched external baselines, independent validation, a prospectively locked seed policy, anatomical quality review, and externally assessed calibration and clinical operating points.
